\documentclass[10pt,a4paper]{amsart}
\usepackage[margin=19mm]{geometry}
\usepackage{amsmath,amssymb,mathtools}
\usepackage[scr=rsfs,cal=euler]{mathalfa}
\usepackage{xfrac}
\usepackage[unicode,hidelinks,bookmarks=false]{hyperref}
\newcommand{\ie}{i.e.\ }
\newcommand{\cf}{cf.\ }
\newcommand{\resp}{respectively\ }
\newcommand{\Subsection}{\subsection}
\providecommand{\nobreakdash}{\nobreak\hbox{-}}
\setlength{\emergencystretch}{2em}
\multlinegap0pt
\usepackage{bm}
\usepackage{bbm}
\usepackage{dsfont}
\usepackage{xspace}
\usepackage{cancel}
\usepackage{mathtools}
\usepackage{xcolor}
\usepackage{amsthm}
\makeatletter
\@ifundefined{lemma}{\newtheorem{lemma}{Lemma}}{}
\@ifundefined{proposition}{\newtheorem{proposition}{Proposition}}{}
\@ifundefined{theorem}{\newtheorem{theorem}{Theorem}}{}
\makeatother
\title[Non-radiating elastic sources in inhomogeneous elastic media]{Non-radiating elastic sources in inhomogeneous elastic media at corners with applications}
\author{Huaian Diao, Yueran Geng, Ruixiang Tang}
\date{Source: arXiv:2502.14468v2 (CC BY 4.0)}
\begin{document}
\maketitle

\noindent\emph{Source and licence.} Non-radiating elastic sources in inhomogeneous elastic media at corners with applications. Version arXiv:2502.14468v2 (CC BY 4.0), distributed under the Creative Commons Attribution 4.0 International licence, as recorded in the arXiv licence metadata for this exact version and shown on the abstract page https://arxiv.org/abs/2502.14468v2. Full rights notice: licenses/arxiv-2502.14468-CC-BY-4.0.txt.\par\smallskip

\noindent\textit{Editorial scope.} Counted unit: the Theorem whose source label is \texttt{thm\_inverse\_problem\_Sh\_f0=0\_nablaf0=0} in the article (source0.tex lines 559--585, source label \texttt{thm\_inverse\_problem\_Sh\_f0=0\_nablaf0=0}), delivered together with the article's own complete original proof of it (source0.tex lines 587--654, one \texttt{proof} environment of 68 lines). No mathematics is written, repaired or re-derived: statement and proof are the article's own text line for line. Required context retained uncounted: ineq\_introduction\_strong\_convexity\_condition (lines 102--104); eq\_introduction\_kpks (lines 120--122); def\_inverse\_problem\_local\_sector\_Sh\_Gammah\_Lambdah (lines 194--196); def\_thm\_2\_ABCD (lines 234--243); lem\_CGOs (lines 359--377); def\_IP\_zeta\_eta (lines 361--363); def\_CGOs\_u0 (lines 365--367); model\_part\_v\_Sh (lines 387--392); eq\_fu0=u0Tnuv\_original (lines 396--398); eq\_IP\_fexpansion\_Bettiformula (lines 404--407); prop\_inverse\_problem\_integral\_re (lines 412--418); ineq\_IP\_dcodtx (lines 421--423); lem\_IP\_est\_int\_Lamdah\_u0Tv-vTu0 (lines 426--432); def\_IP\_u01u02 (lines 441--443); lem\_IP\_est\_f0\_int\_Sh\_u0 (lines 503--510); eq\_IP\_ddperp\_sincosphi (lines 518--527); eq\_IP\_eta\_Taylorexpansion\_e1\_-i1 (lines 529--535). Every definition, display and label of those lines is retained unchanged. Omitted from this package and not redistributed: 1--101, 105--119, 123--193, 197--233, 244--358, 364--364, 368--386, 393--395, 399--403, 408--411, 419--420, 424--425, 433--440, 444--502, 511--517, 528--528, 536--558, 586--586, 655--1291. Nothing inside a retained excerpt is omitted or altered: every retained body is delivered exactly as the article's lines are written, so no omission receipt of any kind accompanies this package. Citations: the retained excerpts carry no citation command at all, so no bibliography entry of the article is transcribed and no reference is redistributed. Reference adaptation: none was needed and none was made; every reference inside the delivered unit resolves inside it, so no unresolved reference or citation is produced. Printed slips kept literal: no printed word, symbol, hypothesis or number of the counted statement or of its proof was altered. Layout and class: the article's own class is \texttt{amsart} with packages including graphicx, amssymb, epstopdf, geometry, array, hyperref, bm, amsmath, amsfonts, amsthm, mathrsfs, cite, color, xcolor; the wrapper is the lane's pinned \texttt{amsart} editorial wrapper at 10pt on A4, which restates the theorem-like environments the retained text needs and adds the article's own macro definitions that the retained text needs (none), with extra wrapper packages (bm, bbm, dsfont, xspace, cancel, mathtools, xcolor). The delivered numbering is the unit's own. Assets: no retained span contains a figure environment, a graphics-inclusion command, a PDF-page inclusion, a table or a TikZ picture; no image file is copied into this package, the package declares no asset dependency and no image file is redistributed with it. Licence: version arXiv:2502.14468v2 of this article is distributed under the Creative Commons Attribution 4.0 International licence, as recorded in the arXiv licence metadata for this exact version and shown on the abstract page https://arxiv.org/abs/2502.14468v2. A full rights notice is delivered at licenses/arxiv-2502.14468-CC-BY-4.0.txt. Characters: every retained excerpt is pure ASCII, so the delivered engine, XeLaTeX with its default Latin Modern text font, reports no missing character.
\begin{align} \label{ineq_introduction_strong_convexity_condition}
	\mu > 0, \quad \lambda + \mu > 0.
\end{align}

\begin{align} \label{eq_introduction_kpks} 
	k_{\mathrm{p}} = \frac{\omega}{\sqrt{\lambda + 2\mu}}, \quad k_{\mathrm{s}} = \frac{\omega}{\sqrt{\mu}}.
\end{align}

\begin{align} \label{def_inverse_problem_local_sector_Sh_Gammah_Lambdah}
	S_{\mathbf{x}_0,h} = W_{\mathbf{x}_0} \cap D_{\mathbf{x}_0,h}, \quad \Gamma_{\mathbf{x}_0,h}^\pm = \Gamma^\pm_{\mathbf{x}_0} \cap D_{\mathbf{x}_0,h}, \quad \Lambda_{\mathbf{x}_0,h} = W_{\mathbf{x}_0} \cap \partial D_{\mathbf{x}_0,h} .
\end{align}  

	\begin{align} \label{def_thm_2_ABCD}
		\begin{cases}
			A=&\sin\left(\theta_M+\theta_m\right)\sin\left(\theta_M-\theta_m\right)\left(1+2\cos\left(\theta_M+\theta_m\right)\cos\left(\theta_M-\theta_m\right)\right), \\
			B=&\sin\left(\theta_M-\theta_m\right)\left(-2\cos\left(\theta_M-\theta_m\right)\cos^2\left(\theta_M+\theta_m\right)+\cos\left(\theta_M-\theta_m\right)\right. \\
			&\left.+\cos\left(\theta_M+\theta_m\right)\right), \\
			C=&\sin\left(\theta_M-\theta_m\right)\left(2\cos\left(\theta_M-\theta_m\right)\cos^2\left(\theta_M+\theta_m\right)-\cos\left(\theta_M-\theta_m\right)\right. \\
			&\left.+\cos\left(\theta_M+\theta_m\right)\right), \\
			D=&\sin\left(\theta_M+\theta_m\right)\sin\left(\theta_M-\theta_m\right)\left(-1+2\cos\left(\theta_M+\theta_m\right)\cos\left(\theta_M-\theta_m\right)\right).
		\end{cases}
	\end{align}

\begin{lemma} \label{lem_CGOs}
	Suppose that the positive function $ \rho\left(\mathbf{x}\right)\in C^2\left(\mathbb{R}^2\right) $ satisfies $ \rho\left(\mathbf{x}\right)=1 $ for $ \mathbf{x}\in\mathbb{R}^2 \backslash \overline{D_R} $. Let $ 0 < R < R_1 < R^{\prime} $, and $ D_{R_1}:=\left\{\mathbf{x} \in \mathbb{R}^2:\left| \mathbf{x} \right|< R_1\right\} $. Recall that $ k_{\mathrm{s}} $ is defined in \eqref{eq_introduction_kpks}. \textcolor{black}{Let $ \mathbf{d} $ and $ \mathbf{d}^{\perp}\in \mathbb{S}^1 $ be} the unit vectors fulfilling that $ \mathbf{d}^{\perp} $ is perpendicular to $ \mathbf{d} $. Assume that 
	\begin{align} \label{def_IP_zeta_eta}
		\bm{\zeta} := \tau \mathbf{d} + \mathrm{i}\sqrt{ \tau^2 + k_{\mathrm{s}}^2 } \mathbf{d}^{\perp}, \quad \bm{\eta} := - \mathrm{i} \sqrt{ 1 + \frac{k_{\mathrm{s}}^2}{\tau^2}} \mathbf{d} + \mathbf{d}^\perp,
	\end{align}
	where $ \tau $ is a constant. If $\tau$ is large enough, then there exists a function  $ \mathbf{R}\left(\mathbf{x}\right)\in C^2\left(D_{R_1}\right) $ such that the Complex Geometrical Optics solution given by
	\begin{align} \label{def_CGOs_u0}
		\mathbf{u}_0 \left(\mathbf{x}\right) := e^{\bm{\zeta} \cdot \mathbf{x}} \bm{\eta} + e^{\bm{\zeta} \cdot \mathbf{x}} \mathbf{R}\left(\mathbf{x}\right), 
	\end{align}
	satisfies the Navier equation 
	\begin{align} \label{eq_CGOs_Navier_equation_no_source}
		\mathcal{L} \mathbf{u} + \omega^2 \rho\left(\mathbf{x}\right) \mathbf{u} = \bm{0} \quad \text{ in } D_{R_1}.
	\end{align}
	Furthermore, when $ \tau $ is sufficiently large, the following estimations hold
	\begin{align} \label{est_inverse_problem_CGO}
		\left\| \mathbf{R} \left(\mathbf{x}\right) \right\|_{L^2 \left(D_{R_1}\right)} \leqslant \frac{c}{\tau}, \quad \left\| \nabla \mathbf{R} \left(\mathbf{x}\right) \right\|_{L^2 \left(D_{R_1}\right)} \leqslant c, \quad \left\| \nabla^2 \mathbf{R} \left(\mathbf{x}\right) \right\|_{L^2 \left(D_{R_1}\right)} \leqslant c \tau, 
	\end{align}
	where $ c $ is a constant depending only on $ R_1 $, $ R^{\prime} $, $ \lambda $, $ \mu $, $ \omega $, and $ \rho\left(\mathbf{x}\right) $.
\end{lemma}

\begin{align} \label{model_part_v_Sh}
	\begin{cases}
		\mathcal{L} \mathbf{v} + \omega^2 \rho\left(\mathbf{x}\right) \mathbf{v} = \mathbf{f} & \text{in } S_{\mathbf{x}_0,h} , \\
		\mathbf{v} = \bm{0}, \ T_{\bm{\nu}} \mathbf{v} = \bm{0} & \text {on } \Gamma_{\mathbf{x}_0,h}^\pm,
	\end{cases}
\end{align}

\begin{align} \label{eq_fu0=u0Tnuv_original}
	\int_{S_h} \mathbf{f} \cdot \mathbf{u}_0 \mathrm{d}\mathbf{x} = \int_{S_h} \mathbf{u}_0 \cdot \mathcal{L} \mathbf{v} - \mathbf{v} \cdot \mathcal{L} \mathbf{u}_0 \mathrm{d}\mathbf{x} = \int_{\Lambda_h} \mathbf{u}_0 \cdot T_{\bm{\nu}} \mathbf{v} - \mathbf{v} \cdot T_{\bm{\nu}} \mathbf{u}_0 \mathrm{d}\sigma.
\end{align}

\begin{align} \label{eq_IP_fexpansion_Bettiformula}
	\mathbf{f} \left(\bm{0}\right) \cdot \int_{S_h} \mathbf{u}_0 \mathrm{d}\mathbf{x} = & \int_{\Lambda_h} \mathbf{u}_0 \cdot T_{\bm{\nu}} \mathbf{v} - \mathbf{v} \cdot T_{\bm{\nu}} \mathbf{u}_0 \mathrm{d}\sigma - \int_{S_h} \left( \nabla \mathbf{f} \left(\bm{0}\right) \mathbf{x} \right) \cdot \mathbf{u}_0 \mathrm{d}\mathbf{x} \notag \\
	& - \int_{S_h} \bm{\delta} \mathbf{f} \left(\mathbf{x}\right) \cdot \mathbf{u}_0 \mathrm{d}\mathbf{x}.
\end{align}

\begin{proposition} \label{prop_inverse_problem_integral_re}
	For positive constants $ \alpha $ and $ \varepsilon $, as long as $ \mathrm{Re}  \left(\mu\right) \geqslant \frac{2\alpha}{e} $, we have
	\begin{align*} 
		\int_{0}^{\varepsilon} r^\alpha e^{- \mu r} \mathrm{d}r = \frac{\Gamma \left( \alpha + 1 \right)}{\mu^{\alpha + 1}} - I_R,
	\end{align*}
	where $ \Gamma \left(s\right)$ denotes the Gamma function, $ I_R = \int_{\varepsilon}^{+\infty} r^\alpha e^{-\mu r} \mathrm{d}r $, $ \left| I_R \right| \leqslant \frac{2}{\mathrm{Re}\left(\mu\right)} e^{ - \frac{\varepsilon}{2} \mathrm{Re} \left(\mu\right)} $. 
\end{proposition}

\begin{align} \label{ineq_IP_dcodtx}
	-1 \leqslant \mathbf{d} \cdot \hat{\mathbf{x}} \leqslant - \delta ,\quad \delta \in \left(0,1\right),
\end{align}

\begin{lemma} \label{lem_IP_est_int_Lamdah_u0Tv-vTu0} 
	Let $ S_h $ and $ \Lambda_h $ be defined in \eqref{def_inverse_problem_local_sector_Sh_Gammah_Lambdah}. Suppose that $ S_h\subset D_R $. Let $ \mathbf{u}_0\left(\mathbf{x}\right) $ be defined in \eqref{def_CGOs_u0} with the unit vector $ \mathbf{d} $ satisfying \eqref{ineq_IP_dcodtx}. Then for $ \mathbf{v}\left(\mathbf{x}\right) \in H^1\left(\mathbb{R}^2\right) $, as $ \tau $ is large enough, it holds that
	\begin{align} \label{est_IP_int_Lamdah_u0Tv-vTu0} 
		\left| \int_{\Lambda_h} \left(\mathbf{u}_0 \cdot \left( T_{\bm{\nu}} \mathbf{v} \right) - \mathbf{v} \cdot \left(T_{\bm{\nu}} \mathbf{u}_0\right)\right) \mathrm{d}\sigma \right| \leqslant c\tau e^{- \delta h \tau},
	\end{align}  
	where the constant $ c $ is independent of $ \tau $.
\end{lemma}

	\textcolor{black}{\begin{align} \label{def_IP_u01u02}
			\mathbf{u}_{01}\left(\mathbf{x}\right) := e^{\bm{\zeta} \cdot \mathbf{x}} \bm{\eta}, \quad \mathbf{u}_{02}\left(\mathbf{x}\right) := e^{\bm{\zeta} \cdot \mathbf{x}} \mathbf{R}\left(\mathbf{x}\right).
	\end{align}}

\begin{lemma} \label{lem_IP_est_f0_int_Sh_u0}
	Consider that $ \mathbf{u}_0\left(\mathbf{x}\right) $ is the CGO solution defined in \eqref{def_CGOs_u0}. Let the convex sectorial corner $ S_h\subset D_R $ be defined as in \eqref{def_inverse_problem_local_sector_Sh_Gammah_Lambdah}. Suppose that the unit vector $ \mathbf{d} $ appearing in \eqref{def_IP_zeta_eta} satisfies \eqref{ineq_IP_dcodtx}. Then for $ \mathbf{f}\left(\bm{0}\right) \in \mathbb{R}^2 $, the estimate
	\begin{align*} 
		\left| \mathbf{f}\left(\bm{0}\right) \cdot \int_{S_h} \mathbf{u}_0\left(\mathbf{x}\right) \mathrm{d}\mathbf{x} \right| \geqslant & \left| \mathbf{f}\left(\bm{0}\right) \right| \left[\left(\theta_M-\theta_m\right)\left(1-\mathcal{O}\left(\frac{1}{\tau^2}\right)\right) \right. \\
		& \left.\left(\frac{1}{3 \tau^2}\left|\mathbf{d}\cdot\hat{\mathbf{x}} \left(\theta_{\xi_1}\right)\right| \left|\mathbf{d}^{\perp}\cdot\hat{\mathbf{x}} \left(\theta_{\xi_1}\right)\right| - \frac{2}{\delta\tau} e^{-\frac{h}{2} \delta \tau}\right) - \frac{c}{\tau}e^{-r_{\varepsilon} \delta \tau}\right] 
	\end{align*}
	holds. Here, $ \tau $ is sufficiently large, $ \hat{\mathbf{x}} \in \Lambda_1 $, $ \theta_{\xi_1} \in \left(\theta_m, \theta_M\right) $, $ \mathbf{d}^{\perp}\cdot\hat{\mathbf{x}} \left(\theta_{\xi_1}\right) \neq 0 $, $ r_{\varepsilon}\in\left(0,h\right) $, and $ c $ is a constant independent of $ \tau $.
\end{lemma}

	\begin{align} \label{eq_IP_ddperp_sincosphi}
		\mathbf{d}=\begin{pmatrix}
			\cos\theta_d \\
			\sin\theta_d
		\end{pmatrix} ,\quad
		\mathbf{d}^\perp=\begin{pmatrix}
			-\sin\theta_d \\
			\cos\theta_d
		\end{pmatrix}, \quad\theta_d\in\left(0,2\pi\right).
	\end{align}

	\begin{align} \label{eq_IP_eta_Taylorexpansion_e1_-i1}
		\bm{\eta} = - \mathrm{i}  \mathbf{d} + \mathbf{d}^\perp - \mathcal{O}\left(\frac{1}{\tau^2} \mathbf{d}\right) = e^{-\mathrm{i}\theta_d} \mathbf{e}_1 - \mathcal{O}\left(\frac{1}{\tau^2} \mathbf{d}\right),\quad
		\mathbf{e}_1:=\begin{pmatrix}
			- \mathrm{i} \\
			1
		\end{pmatrix} .
	\end{align}

\begin{theorem} \label{thm_inverse_problem_Sh_f0=0_nablaf0=0}
	Let $ \lambda $ and $ \mu $ fulfill \eqref{ineq_introduction_strong_convexity_condition} with $ \omega > 0 $. Assume that the positive function $ \rho\left(\mathbf{x}\right)\in C^2\left(\mathbb{R}^2\right) $ and $ \rho\left(\mathbf{x}\right)=1 $ for $ \mathbf{x}\in\mathbb{R}^2 \backslash D_R $. For the sector $ S_{\mathbf{x}_0,h} $ and its boundaries $ \Gamma_{\mathbf{x}_0,h}^\pm $ defined in \eqref{def_inverse_problem_local_sector_Sh_Gammah_Lambdah}, we suppose that $ S_{\mathbf{x}_0,h}\subset D_R $. \textcolor{black}{Let $ \mathbf{v}\left(\mathbf{x}\right) \in H^1\left(S_{\mathbf{x}_0,h}\right) $ satisfy} \eqref{model_part_v_Sh}. If $\mathbf{f}\left(\mathbf{x}\right) \in C^{\alpha} \left(\overline{S_{\mathbf{x}_0,h}}\right)$ with $\alpha\in (0,1)$,  then \begin{align*} 
\mathbf{f} \left(\mathbf{x}_0\right) = \bm{0}. 
	\end{align*}
If $\mathbf{f}\left(\mathbf{x}\right) \in C^{1,\alpha} \left(\overline{S_{\mathbf{x}_0,h}}\right)$ with $\alpha\in (0,1)$,  then 
	\begin{equation} \label{eq_IP_gradientf}
\mathbf{f} \left(\mathbf{x}_0\right) = \bm{0},\	\begin{aligned}
			\nabla\mathbf{f}\left(\mathbf{x}_0\right)
			\begin{pmatrix}
				A \\ 
				B
			\end{pmatrix}
			+
			\begin{pmatrix}
				0 & -1 \\
				1 & 0
			\end{pmatrix}
			\nabla\mathbf{f}\left(\mathbf{x}_0\right)
			\begin{pmatrix}
				C \\ 
				D
			\end{pmatrix}
			=\mathbf{0},
		\end{aligned}
	\end{equation}
	with $ A $, $ B $, $ C $, and $ D $ defined in \eqref{def_thm_2_ABCD}. 
\end{theorem}

\begin{proof} 
	Without loss of generality, we suppose that $ \mathbf{x}_0 = \bm{0} $, as the operator $ \mathcal{L} $ is invariant under rigid motion. In the following, we deal with the case that $\mathbf{f}\left(\mathbf{x}\right) \in C^{1,\alpha} \left(\overline{S_{\mathbf{0},h}}\right)$. The case for $\mathbf{f}\left(\mathbf{x}\right) \in C^{\alpha} \left(\overline{S_{\mathbf{0},h}}\right)$ can be proved in a similar way. By virtue of Lemma \ref{lem_CGOs} and the first Betti's formula, we have \eqref{eq_fu0=u0Tnuv_original} and \eqref{eq_IP_fexpansion_Bettiformula}. Then it follows that
	\begin{align*} 
		\left| \mathbf{f} \left(\bm{0}\right) \cdot \int_{S_h} \mathbf{u}_0 \mathrm{d}\mathbf{x} \right|	\leqslant & \left| \int_{\Lambda_h} \mathbf{u}_0 \cdot T_{\bm{\nu}} \mathbf{v} - \mathbf{v} \cdot T_{\bm{\nu}} \mathbf{u}_0 \mathrm{d}\sigma \right| + \left| \int_{S_h} \left( \nabla \mathbf{f} \left(\bm{0}\right) \mathbf{x} \right) \cdot \mathbf{u}_0 \mathrm{d}\mathbf{x} \right| \\
		& + \left| \int_{S_h} \bm{\delta} \mathbf{f} \left(\mathbf{x}\right) \cdot \mathbf{u}_0 \mathrm{d}\mathbf{x} \right|. 
	\end{align*}
	With the help of Lemma \ref{lem_IP_est_int_Lamdah_u0Tv-vTu0}-\ref{lem_IP_est_f0_int_Sh_u0}, as $ \tau $ is sufficiently large, we have  
	\begin{align} \label{ineq_inverse_problem_f0_tau^2}
		\left| \mathbf{f}\left(\bm{0}\right) \right|\left[\left(\theta_M-\theta_m\right)\left(1-\mathcal{O}\left(\frac{1}{\tau^2}\right)\right) \left(\frac{1}{3 \tau^2}\left|\mathbf{d}\cdot\hat{\mathbf{x}} \left(\theta_{\xi_1}\right)\right| \left|\mathbf{d}^{\perp}\cdot\hat{\mathbf{x}} \left(\theta_{\xi_1}\right)\right| - \frac{2}{\delta\tau} e^{-\frac{h}{2} \delta \tau}\right) \right. \notag \\
		\left.- \frac{c}{\tau}e^{-r_{\varepsilon} \delta \tau}\right] \leqslant  c \tau e^{- \delta h \tau} +  \frac{c}{\tau^3} + \frac{c}{\tau^{3+\alpha}}.
	\end{align}
	Here, $ c $ is a constant independent of $ \tau $, $ r_{\varepsilon}\in\left(0,h\right) $, $ \theta_{\xi_1} \in \left(\theta_m, \theta_M\right) $, $ \mathbf{d} $, $ \mathbf{d}^{\perp} $ and $ \hat{\mathbf{x}}\left(\theta_{\xi_1}\right) $ are unit vectors satisfying \eqref{ineq_IP_dcodtx} and $ \mathbf{d}^{\perp}\cdot\hat{\mathbf{x}} \left(\theta_{\xi_1}\right) \neq 0 $. Multiplying $ \tau^2 $ on the both sides of \eqref{ineq_inverse_problem_f0_tau^2}, we obtain 
	\begin{align*} 
		&\left(\theta_M-\theta_m\right)\left(\frac{1}{3 }\left|\mathbf{d}\cdot\hat{\mathbf{x}} \left(\theta_{\xi_1}\right)\right| \left|\mathbf{d}^{\perp}\cdot\hat{\mathbf{x}} \left(\theta_{\xi_1}\right)\right| - \frac{2}{\delta}\tau e^{-\frac{h}{2} \delta \tau}\right) \left| \mathbf{f}\left(\bm{0}\right) \right| \\
		\leqslant &\frac{c}{\tau} +  \frac{c}{\tau^{1+\alpha}} + \mathcal{O}\left(\frac{1}{\tau^2}\right) + c\tau e^{-r_{\varepsilon} \delta \tau} + c \tau^3 e^{- \delta h \tau} .
	\end{align*}
	Letting $ \tau\to+\infty $, it yields that $ \mathbf{f} \left(\bm{0}\right) = \bm{0} $.
	
	Let $ \mathbf{d} $, $ \mathbf{d}^{\perp} $ and $ \mathbf{e}_1 $ be formulated as in \eqref{eq_IP_ddperp_sincosphi} and \eqref{eq_IP_eta_Taylorexpansion_e1_-i1} respectively. When $ \tau $ is sufficiently large, $ \bm{\zeta} $ and $ \bm{\eta} $ appearing in Lemma \ref{lem_CGOs} can be formulated as
	\begin{align}
		\bm{\zeta} = \tau \left( \mathbf{d} + \text{i} \mathbf{d}^{\perp} \right) + \mathcal{O} \left(\frac{1}{\tau}\mathbf{d}^{\perp}\right) = \tau \mathrm{i} e^{-\mathrm{i}\theta_d} \mathbf{e}_1  + \mathcal{O} \left(\frac{1}{\tau}\mathbf{d}^{\perp}\right), \notag
	\end{align}
	and \eqref{eq_IP_eta_Taylorexpansion_e1_-i1}, respectively. \textcolor{black}{Recall $\mathbf{u}_{01}$ and $\mathbf{u}_{02}$ defined in \eqref{def_IP_u01u02}.} Then we have
	\begin{align} \label{eq_intSh_fu01}
		& \int_{S_h} \mathbf{f} \cdot \mathbf{u}_{01} \mathrm{d}\mathbf{x} = \int_{S_h} e^{\bm{\zeta} \cdot \mathbf{x}} \left[\left(\nabla \mathbf{f} \left(\bm{0}\right) \mathbf{x}\right)\cdot\bm{\eta} + \bm{\delta} \mathbf{f}\left(\mathbf{x}\right)\cdot \bm{\eta}\right] \mathrm{d}\mathbf{x} \notag \\
		= & \int_{S_h} e^{\left[\tau \mathrm{i} e^{-\mathrm{i}\theta_d} \mathbf{e}_1 \cdot \mathbf{x} + \mathcal{O} \left(\frac{1}{\tau}\mathbf{d}^{\perp} \cdot \mathbf{x}\right)\right]}\left[ \left(\nabla \mathbf{f} \left(\bm{0}\right) \mathbf{x}\right)\cdot\mathbf{e}_1 e^{-\mathrm{i}\theta_d}\right. + \mathcal{O} \left(\frac{1}{\tau^2}\left(\nabla \mathbf{f} \left(\bm{0}\right) \mathbf{x}\right)\cdot\mathbf{d}\right)\notag \\
		& \left. + \mathcal{O}\left(\left|\mathbf{x}\right|^{1+\alpha}\right)\right] \mathrm{d}\mathbf{x} \notag \\
		= & \int_{\theta_m}^{\theta_M} \int_{0}^{h} r^2 e^{r \tau\left[e^{\mathrm{i}\left(\theta-\theta_d\right)} + \mathcal{O}\left(\frac{1}{\tau^2}\right)\right]}\left[e^{-\mathrm{i}\theta_d}g\left(\theta\right) + \mathcal{O}\left(\frac{1}{\tau^2}\right) + \mathcal{O}\left(r^{\alpha}\right)\right] \mathrm{d}r \mathrm{d}\theta ,
	\end{align}
	where $ g\left(\theta\right) := -\text{i}\left(\frac{\partial f_1}{\partial x_1}\left(\bm{0}\right)\cos\theta + \frac{\partial f_1}{\partial x_2}\left(\bm{0}\right)\sin\theta\right) + \left(\frac{\partial f_2}{\partial x_1}\left(\bm{0}\right)\cos\theta + \frac{\partial f_2}{\partial x_2}\left(\bm{0}\right)\sin\theta\right) $. Due to \eqref{ineq_IP_dcodtx}, for all $ \theta\in\left[\theta_m, \theta_M\right] $, it holds that $ \cos \left(\theta-\theta_d\right) \leqslant -\delta < 0 $. By Proposition \ref{prop_inverse_problem_integral_re}, we obtain
	\begin{align*}
		& \int_{\theta_m}^{\theta_M} \int_{0}^{h} r^2 e^{r \tau\left[e^{\mathrm{i}\left(\theta-\theta_d\right)} + \mathcal{O}\left(\frac{1}{\tau^2}\right)\right]} e^{-\mathrm{i}\theta_d}g\left(\theta\right) \mathrm{d}r \mathrm{d}\theta \\
		= & \int_{\theta_m}^{\theta_M} \frac{\Gamma\left(3\right) }{-\tau^3\left[e^{\text{i}\left(\theta-\theta_d\right)} + \mathcal{O} \left(\frac{1}{\tau^2}\right)\right]^3} e^{-\text{i}\theta_d}g\left(\theta\right) \mathrm{d}\theta - \int_{\theta_m}^{\theta_M} e^{-\text{i}\theta_d} I_R\left(\theta\right) g\left(\theta\right) \mathrm{d}\theta \\
		= & \left(\frac{2e^{2\text{i}\theta_d}}{-\tau^3} + \mathcal{O} \left(\frac{1}{\tau^5}\right)\right)  \int_{\theta_m}^{\theta_M} e^{-3\text{i}\theta} g\left(\theta\right) \mathrm{d}\theta + \mathcal{O} \left( \frac{1}{\tau} e^{-\frac{h}{2}\delta\tau} \right) ,
	\end{align*}
	where $ I_R\left(\theta\right) := \int_{h}^{+\infty} r^2 e^{r \tau\left[e^{\mathrm{i}\left(\theta-\theta_d\right)} + \mathcal{O}\left(\frac{1}{\tau^2}\right)\right]} \mathrm{d}r $, and 
	\begin{align*}
		\left| I_R\left(\theta\right) \right| \leqslant \frac{2}{-\tau\left[\cos\left(\theta-\theta_d\right) + \mathcal{O}\left(\frac{1}{\tau^2}\right)\right]}e^{\frac{h}{2}\tau\left[\cos\left(\theta-\theta_d\right) + \mathcal{O}\left(\frac{1}{\tau^2}\right)\right]} .
	\end{align*}
	Similarly, it can be calculated that
	\begin{align*}
		\int_{\theta_m}^{\theta_M} \int_{0}^{h} r^2 e^{r \tau\left[e^{\mathrm{i}\left(\theta-\theta_d\right)} + \mathcal{O}\left(\frac{1}{\tau^2}\right)\right]} \mathcal{O}\left(\frac{1}{\tau^2}\right) \mathrm{d}r \mathrm{d}\theta = \mathcal{O} \left(\frac{1}{\tau^5}\right) ,\\
		\int_{\theta_m}^{\theta_M} \int_{0}^{h} r^2 e^{r \tau\left[e^{\mathrm{i}\left(\theta-\theta_d\right)} + \mathcal{O}\left(\frac{1}{\tau^2}\right)\right]} \mathcal{O}\left(r^{\alpha}\right) \mathrm{d}r \mathrm{d}\theta = \mathcal{O} \left(\frac{1}{\tau^{3+\alpha}}\right) .
	\end{align*}
	Then, \eqref{eq_intSh_fu01} can be rewritten as
	\begin{align} \label{eq_IP_int_Sh_fu01}
		\int_{S_h} \mathbf{f} \cdot \mathbf{u}_{01} \mathrm{d}\mathbf{x} = & \left(\frac{2e^{2\text{i}\theta_d}}{-\tau^3} + \mathcal{O} \left(\frac{1}{\tau^5}\right)\right)  \int_{\theta_m}^{\theta_M} e^{-3\text{i}\theta} g\left(\theta\right) \mathrm{d}\theta + \mathcal{O} \left( \frac{1}{\tau} e^{-\frac{h}{2}\delta\tau} \right) + \mathcal{O} \left(\frac{1}{\tau^5}\right)\notag \\
		&  + \mathcal{O} \left(\frac{1}{\tau^{3+\alpha}}\right) .
	\end{align}
	By \eqref{eq_fu0=u0Tnuv_original}, Lemma \ref{lem_IP_est_int_Lamdah_u0Tv-vTu0}, and integral mean value theorem, as $ \tau $ is large enough, we arrive that
	\begin{align} \label{ineq_IP_int_Sh_fu01}
		\left|\int_{S_h} \mathbf{f} \cdot \mathbf{u}_{01} \mathrm{d}\mathbf{x}\right| & \leqslant \left|\int_{S_h} \mathbf{f} \cdot \mathbf{u}_{0} \mathrm{d}\mathbf{x}\right| + \left|\int_{S_h} \mathbf{f} \cdot \mathbf{u}_{02} \mathrm{d}\mathbf{x}\right| \notag \\
		& \leqslant \left|\int_{\Lambda_h} \mathbf{u}_0 \cdot T_{\bm{\nu}} \mathbf{v} - \mathbf{v} \cdot T_{\bm{\nu}} \mathbf{u}_0 \mathrm{d}\sigma\right| + ce^{-\delta\left|\mathbf{x}_{\varepsilon}\right|\tau} \left|\mathbf{f}\left(\mathbf{x}_{\varepsilon}\right)\right| \left\|\mathbf{R}\left(\mathbf{x}\right)\right\|_{L^2\left(S_h\right)} \notag \\
		& \leqslant c\tau e^{- \delta h \tau} +  \frac{c}{\tau}e^{-\delta\left|\mathbf{x}_{\varepsilon}\right|\tau} ,
	\end{align}
	where $ \mathbf{x}_{\varepsilon} \in S_h $, $ \left|\mathbf{x}_{\varepsilon}\right|\in\left(0,h\right) $, and $ c $ is a constant independent of $ \tau $. Combining \eqref{eq_IP_int_Sh_fu01} and \eqref{ineq_IP_int_Sh_fu01}, it can be formulated that
	\begin{align} \label{eq_IP_gtau^3}
		& \left|\frac{2e^{2\text{i}\theta_d}}{-\tau^3}  \int_{\theta_m}^{\theta_M} e^{-3\text{i}\theta} g\left(\theta\right) \mathrm{d}\theta + \mathcal{O} \left( \frac{1}{\tau} e^{-\frac{h}{2}\delta\tau} \right) + \mathcal{O} \left(\frac{1}{\tau^5}\right) + \mathcal{O} \left(\frac{1}{\tau^{3+\alpha}}\right)\right| \notag \\
		\leqslant & c\tau e^{- \delta h \tau} + \frac{c}{\tau}e^{-\delta\left|\mathbf{x}_{\varepsilon}\right|\tau} .
	\end{align}
	Multiply $ \tau^3 $ on the both ends of \eqref{eq_IP_gtau^3} and let $ \tau \to +\infty $. It yields that
	\begin{align} \label{eq_inverse_problem_int_eg=0}
		\int_{\theta_m}^{\theta_M} e^{-3\text{i}\theta} g\left(\theta\right) \mathrm{d}\theta = 0 . 
	\end{align}
	From the real and imaginary parts of \eqref{eq_inverse_problem_int_eg=0}, we arrive at \eqref{eq_IP_gradientf} with $ A $, $ B $, $ C $, and $ D $ defined in \eqref{def_thm_2_ABCD}.
	
	The proof is complete.
\end{proof}

\end{document}
